\documentclass[11pt]{article}
% ===================================================================
% Shared preamble for the "tiny lemma, read slowly" preprint series.
% Mirrors the style of FrobeniusLemmas_preprint.tex.
% ===================================================================

\usepackage{amsmath,amssymb,amsthm}
\usepackage{mathtools}
\usepackage[margin=1.2in]{geometry}
\usepackage{hyperref}
\usepackage{xcolor}
\usepackage{listings}
\usepackage{tikz}
\usepackage{tikz-cd}
\usepackage{mathpartir}
\usepackage{newunicodechar}
\usepackage{setspace}

\onehalfspacing

% ---- Unicode glyphs ----
\newunicodechar{→}{\ensuremath{\to}}
\newunicodechar{↦}{\ensuremath{\mapsto}}
\newunicodechar{↔}{\ensuremath{\leftrightarrow}}
\newunicodechar{⟹}{\ensuremath{\Longrightarrow}}
\newunicodechar{⟶}{\ensuremath{\longrightarrow}}
\newunicodechar{≠}{\ensuremath{\neq}}
\newunicodechar{≤}{\ensuremath{\leq}}
\newunicodechar{≥}{\ensuremath{\geq}}
\newunicodechar{∀}{\ensuremath{\forall}}
\newunicodechar{∃}{\ensuremath{\exists}}
\newunicodechar{∘}{\ensuremath{\circ}}
\newunicodechar{·}{\ensuremath{\cdot}}
\newunicodechar{∈}{\ensuremath{\in}}
\newunicodechar{∉}{\ensuremath{\notin}}
\newunicodechar{≃}{\ensuremath{\simeq}}
\newunicodechar{≅}{\ensuremath{\cong}}
\newunicodechar{≈}{\ensuremath{\approx}}
\newunicodechar{∧}{\ensuremath{\wedge}}
\newunicodechar{∨}{\ensuremath{\vee}}
\newunicodechar{¬}{\ensuremath{\neg}}
\newunicodechar{⊢}{\ensuremath{\vdash}}
\newunicodechar{⟨}{\ensuremath{\langle}}
\newunicodechar{⟩}{\ensuremath{\rangle}}
\newunicodechar{⇑}{\ensuremath{\mathord{\uparrow}}}
\newunicodechar{↑}{\ensuremath{\mathord{\uparrow}}}
\newunicodechar{∎}{\ensuremath{\blacksquare}}
\newunicodechar{∣}{\ensuremath{\mid}}
\newunicodechar{∤}{\ensuremath{\nmid}}
\newunicodechar{∑}{\ensuremath{\textstyle\sum}}
\newunicodechar{∏}{\ensuremath{\textstyle\prod}}
\newunicodechar{∅}{\ensuremath{\emptyset}}
\newunicodechar{∪}{\ensuremath{\cup}}
\newunicodechar{∩}{\ensuremath{\cap}}
\newunicodechar{≡}{\ensuremath{\equiv}}
\newunicodechar{ℕ}{\ensuremath{\mathbb{N}}}
\newunicodechar{ℤ}{\ensuremath{\mathbb{Z}}}
\newunicodechar{ℝ}{\ensuremath{\mathbb{R}}}
\newunicodechar{𝔽}{\ensuremath{\mathbb{F}}}
\newunicodechar{α}{\ensuremath{\alpha}}
\newunicodechar{β}{\ensuremath{\beta}}
\newunicodechar{σ}{\ensuremath{\sigma}}
\newunicodechar{τ}{\ensuremath{\tau}}
\newunicodechar{φ}{\ensuremath{\varphi}}
\newunicodechar{ψ}{\ensuremath{\psi}}
\newunicodechar{χ}{\ensuremath{\chi}}
\newunicodechar{Φ}{\ensuremath{\Phi}}
\newunicodechar{Δ}{\ensuremath{\Delta}}
\newunicodechar{₀}{\textsubscript{0}}
\newunicodechar{₁}{\textsubscript{1}}
\newunicodechar{₂}{\textsubscript{2}}
\newunicodechar{ᵏ}{\ensuremath{{}^{k}}}
\newunicodechar{ⁿ}{\ensuremath{{}^{n}}}
\newunicodechar{ʲ}{\ensuremath{{}^{j}}}
\newunicodechar{²}{\ensuremath{{}^{2}}}
\newunicodechar{³}{\ensuremath{{}^{3}}}
\newunicodechar{⁴}{\ensuremath{{}^{4}}}
\newunicodechar{⁻}{\ensuremath{{}^{-}}}
\newunicodechar{¹}{\ensuremath{{}^{1}}}

% ---- Lean listing style ----
\definecolor{leanbg}{RGB}{247,247,250}
\definecolor{leankw}{RGB}{0,90,160}
\definecolor{leancm}{RGB}{120,120,120}
\lstdefinestyle{lean}{
  backgroundcolor=\color{leanbg},
  basicstyle=\ttfamily\small,
  keywordstyle=\color{leankw}\bfseries,
  commentstyle=\color{leancm}\itshape,
  morekeywords={theorem,lemma,def,fun,Prop,Type,variable,where,
    Field,Fintype,Finite,DecidableEq,CommRing,CommSemiring,Ring,CharP,Function,
    Injective,Surjective,Bijective,by,rfl,have,rw,exact,ext,simp,intro,
    iterateFrobenius,RingHom,Commute,convert,apply,using,Nat,Coprime,gcd,
    Odd,Even,IsAPN,IsAB,walsh,Nat,obtain,induction,cases,calc,grind},
  showstringspaces=false,
  columns=fullflexible,
  extendedchars=true,
  frame=single,
  framesep=6pt,
  xleftmargin=4pt,
  aboveskip=12pt,
  belowskip=14pt,
}
\lstset{style=lean}

\newtheorem{lemma}{Lemma}
\newtheorem{theorem}{Theorem}
\newtheorem{corollary}{Corollary}

% boxed "what just happened" note
\newcommand{\note}[1]{%
  \par\smallskip
  \noindent\colorbox{leanbg}{\parbox{\dimexpr\linewidth-2\fboxsep}{\small #1}}%
  \par\smallskip}

\usepackage{colortbl}

% ---- a few extra glyphs used only in this note ----
\newunicodechar{κ}{\ensuremath{\kappa}}
\newunicodechar{ᵢ}{\ensuremath{{}_{i}}}
\newunicodechar{ₘ}{\ensuremath{{}_{m}}}
\newunicodechar{ₐ}{\ensuremath{{}_{a}}}
\newunicodechar{₃}{\textsubscript{3}}
\newunicodechar{Σ}{\ensuremath{\textstyle\sum}}
\newunicodechar{∥}{\ensuremath{\Vert}}

\title{\textbf{One sorry, named}\\[4pt]
\large The single missing node of NSUCRYPTO'2019 Problem~11, read slowly}
\author{}
\date{}

\begin{document}
\maketitle

\begin{abstract}
\noindent
This is a \emph{gap report}, not a proof.\\[4pt]
The whole derivation of NSUCRYPTO'2019 Problem~11 (and the unconditional
$m$-tuple count for the Kasami family) has been pushed down to \emph{one}
self-standing statement.  Every other layer --- the concrete character, the
two framework bridges, the autocorrelation homogeneity, the assembly into the
closed-form count --- is machine-checked in Lean~4 with no \texttt{sorry} and
only the standard axioms.\\[6pt]
\noindent
What remains is a single harmonic-analytic fact, the
\emph{almost-bent sum-zero correlation balance}:
\[
  \sum_{v\neq 0}\ \prod_{i} \mathrm{Wcorr}\bigl(\chi,\, x^{d_k},\, 1,\, v\,c_i\bigr)
  \;=\; 0
  \qquad\Bigl(\textstyle\sum_i c_i = 0,\ c_i\neq 0\Bigr).
\]
\noindent
We state \emph{exactly} what is proved, \emph{exactly} where the hole is, read
it as \emph{math} $\leftrightarrow$ \emph{Lean} $\leftrightarrow$ a missing
piece of \emph{spectral infrastructure}, and explain why the obvious inputs
(three-valued Walsh spectrum + homogeneity) come up one structural step short.
\end{abstract}

\bigskip

% =====================================================================
\section*{The shape of the gap}

\noindent
Almost everything is a staircase that ends one step below the landing.
The single open node sits between two fully-proved layers:

\[
\begin{tikzcd}[row sep=large, column sep=small]
\boxed{\texttt{kasami\_is\_ab}}
  \arrow[d, Rightarrow, "\text{(proved: 3-valued spectrum)}"'] &[-2.2em] \\
\textbf{[KERNEL]}\ \ \texttt{kasami\_correlationBalance\_sumzero}
  \arrow[d, Rightarrow, "\text{(proved: }\texttt{correlationBalance\_iff\_count}\text{)}"'] \\
\texttt{kasami\_triple\_count\_unconditional}\ \ (\kappa_3 = 2^{2n-3})
  \arrow[d, Rightarrow] \\
\texttt{nsucrypto\_problem11}
\end{tikzcd}
\]

\note{The box at the top is a \emph{theorem} (Kasami's Walsh spectrum is
three-valued, proved with standard axioms only).  The node in \textbf{bold} is
the \emph{only} \texttt{sorry} in the entire chain.  Everything below it is a
term that consumes the kernel and produces the closed-form counts.}

\bigskip

% =====================================================================
\section{What is solid below the gap}

\noindent
Fix a finite field $\mathbb{F}=\mathrm{GF}(2^n)$.  An \emph{additive character}
is packaged as a structure \texttt{Chi}; the project supplies a \emph{concrete}
one, the trace sign character $\chi(x)=(-1)^{\mathrm{Tr}(x)}$:

\begin{lstlisting}
noncomputable def traceChi : MTupleCount.Chi 𝔽 where
  app      := WalshAB.χ
  app_zero := WalshAB.χ_zero
  app_add  := WalshAB.χ_mul
  orth     := WalshAB.χ_sum_eq
\end{lstlisting}

\note{This discharges the \emph{existence} of the abstract character the whole
spectral machinery was parameterised over.  Proved, no \texttt{sorry}.}

\medskip

\noindent
The \textbf{derivative autocorrelation} at shift $a$ and mask $w$ is
\[
  \mathrm{Wcorr}(\chi, f, a, w)
  \;=\;
  \sum_{x\in\mathbb{F}} \chi\!\bigl(w\cdot(f(x+a)+f(x))\bigr)\in\mathbb{Z},
\]
and a single proved lemma makes the shift $a$ irrelevant for a power map:

\begin{lstlisting}
lemma wcorr_pow_homog (χ : MTupleCount.Chi 𝔽) (d : ℕ) (hd : 0 < d) (a w : 𝔽) :
    MTupleCount.Wcorr χ (fun x => x ^ d) a w
      = MTupleCount.Wcorr χ (fun x => x ^ d) 1 (w * a ^ d)
\end{lstlisting}

\note{\emph{Homogeneity.}  Every derivative direction $a$ is a relabelling of
$a=1$.  This is why the kernel below may fix $a=1$ with no loss.  Proved, no
\texttt{sorry}.}

\medskip

\noindent
Two \textbf{framework bridges} identify the two counting worlds used in the
project (collision counting vs.\ Fourier counting), so the conjecture's
pure-counting statement can actually be fed to the spectral machinery:

\begin{lstlisting}
lemma Delta0_eq_Δ   (f : 𝔽 → 𝔽) (a : 𝔽) :
    M002.derivImage f a = MTupleCount.Δ f a
lemma tcount_eq_kappa (S : Finset 𝔽) (v₁ v₂ : 𝔽) :
    tcount S v₁ v₂ = MTupleCount.κ 3 S ![v₁, v₂, v₁ + v₂]
\end{lstlisting}

\note{Both proved, no \texttt{sorry}.  In characteristic two $-$ and $+$ agree,
so the raw derivative image is literally $\Delta f\,a$.}

\bigskip

% =====================================================================
\section{What is solid above the gap}

\noindent
The bridge from the kernel to the count is a proved \emph{equivalence}, not a
guess.  Correlation balance is exactly the counting identity:

\[
  \underbrace{\sum_{v\neq 0}\ \prod_i \mathrm{Wcorr}(\chi,f,a,v\,c_i)=0}
            _{\textstyle \texttt{CorrelationBalance}}
  \quad\Longleftrightarrow\quad
  |\mathbb{F}|\cdot \kappa_m \;=\; |\Delta f\,a|^{\,m}.
\]

\begin{lstlisting}
lemma correlationBalance_iff_count (χ : Chi 𝔽) (f : 𝔽 → 𝔽) (hf : APN f)
    (a : 𝔽) (ha : a ≠ 0) (m : ℕ) (c : Fin m → 𝔽) :
    CorrelationBalance χ f a m c ↔
      (card 𝔽 : ℤ) * (κ m (Δ f a) c) = ((Δ f a).card : ℤ) ^ m
\end{lstlisting}

\note{Proved, no \texttt{sorry}.  Each product is $2^m\,P(v)$ by the APN
fibre-doubling bridge; summing over $v\neq0$ and peeling the $v=0$ term gives
$2^m\bigl(|\mathbb{F}|\,\kappa_m - |\Delta|^m\bigr)$, which vanishes iff the
counting identity holds.}

\medskip

\noindent
Hence, \emph{given the kernel}, the closed forms are assembled with no further
hypotheses:

\begin{lstlisting}
theorem kasami_triple_count_unconditional ... :
    MTupleCount.κ 3 (MTupleCount.Δ (fun x : 𝔽 => x ^ KasamiAPN.kasamiExp k) 1)
        ![v₁, v₂, v₁ + v₂] = 2 ^ (2 * n - 3)
\end{lstlisting}

\note{This term \emph{compiles}.  Its only transitive \texttt{sorry} is the one
kernel below.  Remove the kernel's \texttt{sorry} and the whole Problem~11
chain is closed.}

\bigskip

% =====================================================================
\section{The gap itself, stated exactly}

\noindent
Here is the one open statement, verbatim from the source.  The Kasami exponent
is $d_k = 2^{2k}-2^k+1$ (\texttt{kasamiExp k}); the coefficient triple
$\,[\,v_1,\,v_2,\,v_1{+}v_2\,]\,$ \emph{sums to zero} in characteristic two.

\begin{lstlisting}
/-- [KERNEL] Almost-bent sum-zero triple balance for the Kasami map. -/
theorem kasami_correlationBalance_sumzero (n k : ℕ) (hn : 3 ≤ n)
    (hk1 : 1 ≤ k) (hkn : k < n) (hodd : Odd n)
    (hcop : Nat.Coprime k n) (hcard : card 𝔽 = 2 ^ n)
    (v₁ v₂ : 𝔽) (hv₁ : v₁ ≠ 0) (hv₂ : v₂ ≠ 0) (hne : v₁ ≠ v₂) :
    MTupleCount.CorrelationBalance (traceChi (𝔽 := 𝔽))
      (fun x : 𝔽 => x ^ KasamiAPN.kasamiExp k) 1 3 ![v₁, v₂, v₁ + v₂] := by
  sorry          -- ← the entire missing content of Problem 11
\end{lstlisting}

\medskip

\noindent
Unfolded, the goal is the single scalar identity
\[
  \sum_{v\in\mathbb{F}^{\times}}
    \mathrm{Wcorr}(v\,v_1)\;
    \mathrm{Wcorr}(v\,v_2)\;
    \mathrm{Wcorr}\bigl(v\,(v_1{+}v_2)\bigr)
  \;=\;0,
  \qquad
  \mathrm{Wcorr}(w):=\mathrm{Wcorr}\bigl(\chi,\,x^{d_k},\,1,\,w\bigr).
\]

\note{It is one number being zero.  Not a quantifier over fibres, not a
cardinality bound --- a \emph{cancellation}.  That is what makes it the genuine
harmonic-analytic heart of the problem, and what makes it resist the
elementary moves that closed every layer around it.}

\bigskip

% =====================================================================
\section{Why the obvious inputs fall one step short}

\noindent
The temptation is: ``Kasami is almost bent, so its Walsh spectrum is
three-valued; substitute and cancel.''  Here is precisely how far that goes,
and where it stops.

\subsection*{What three-valuedness gives (a single mask)}

\noindent
\texttt{kasami\_is\_ab} says: for each fixed nonzero output mask $w$,
\[
  W(\chi,\,x^{d_k},\,u,\,w)\ \in\ \bigl\{\,0,\ \pm 2^{(n+1)/2}\,\bigr\}
  \qquad\text{for every input mask }u .
\]
By Wiener--Khinchin (proved, \texttt{card\_mul\_corr}) the autocorrelation is the
\emph{inverse} Fourier transform of the squared spectrum,
\[
  |\mathbb{F}|\cdot\mathrm{Wcorr}(\chi,f,a,w)
  \;=\;
  \sum_{u} \chi(u\cdot a)\,\bigl(W(\chi,f,u,w)\bigr)^2 ,
\]
and the second/third moments of a \emph{single} three-valued column are pinned
down (\texttt{second\_moment}, \texttt{third\_moment\_ab}).  So for one mask at
a time the spectral data is completely understood.

\subsection*{What the kernel actually needs (a product of three masks)}

\noindent
The kernel is not about one column.  Expanding each $\mathrm{Wcorr}$ through
Wiener--Khinchin and multiplying, the sum-zero balance becomes a statement about
the \emph{joint} support of the spectrum across the \emph{three different output
masks} $v\,v_1,\ v\,v_2,\ v\,(v_1{+}v_2)$ simultaneously:
\[
  \sum_{v\neq 0}\ \prod_{i=1}^{3}
     \Bigl(\textstyle\sum_{u_i}\chi(u_i)\,W(u_i,\,v\,c_i)^2\Bigr)
  \;=\;0 .
\]
The cancellation is forced by \emph{how the three nonzero values $\pm
2^{(n+1)/2}$ are distributed jointly} --- the cross-correlation /
\emph{joint-support} structure of the Kasami Walsh transform.  Knowing each
factor is three-valued says nothing about the \emph{correlation between the
factors}, and the sum-zero condition $c_1+c_2+c_3=0$ is exactly what couples
them.

\[
\begin{tikzcd}[row sep=small, column sep=large]
\text{3-valued, per mask} \arrow[r, Rightarrow, "\text{known}"]
   & \text{moments of one column} \\
\text{\textbf{joint support, 3 masks}} \arrow[r, Rightarrow, dashed, "\text{\textbf{missing}}"]
   & \text{sum-zero cancellation}
\end{tikzcd}
\]

\note{The proved layer controls each \emph{factor}; the kernel needs the
\emph{joint distribution} of the factors.  That is one quantifier deeper than
\texttt{AlmostBent}, and Mathlib carries no API for the $m$-fold support
geometry of a three-valued spectrum.}

\bigskip

% =====================================================================
\section{The same hole, in four equivalent contexts}

\noindent
Because the surrounding bridges are proved equivalences, the one missing fact
can be phrased in whichever language one prefers; each phrasing is provably the
\emph{same} node, and each is open for the \emph{same} reason.

\begin{center}
\renewcommand{\arraystretch}{1.7}
\begin{tabular}{l l}
\textbf{context} & \textbf{the missing statement} \\
\hline
Walsh correlation & $\displaystyle\sum_{v\neq0}\prod_i \mathrm{Wcorr}(v\,c_i)=0$ \quad (the kernel) \\
Fourier counting  & $|\mathbb{F}|\cdot\kappa_3(c) = |\Delta|^3$ \\
spectral correction & $\mathrm{tripleError}(\chi,\Delta,c)=0$ \\
original collision & $|\mathbb{F}|\cdot\mathrm{tcount}(\Delta_0) = |\Delta_0|^3$ \\
\end{tabular}
\end{center}

\note{Bridges: \texttt{correlationBalance\_iff\_count},
\texttt{kernel\_iff\_constCount}, \texttt{kernel\_iff\_tripleError\_zero},
\texttt{tcount\_eq\_kappa} --- all proved.  Closing the kernel in \emph{any}
one column closes \emph{all} of them.}

\bigskip

% =====================================================================
\section*{What would close it}

\noindent
The remaining content is the Carlet--Charpin--Zinoviev almost-bent
$m$-balance.  In Lean terms the missing infrastructure is the
\emph{joint-support / higher-moment} theory of an almost-bent spectrum:
the value of
\[
  \sum_{w}\ \prod_{i=1}^{m}\bigl(W(u_i,\,w)\bigr)^2
\]
over a three-valued spectrum, with the coupling fixed by $\sum_i c_i = 0$.
This is the one piece not currently in Mathlib, and the one piece this note
flags as open.

\bigskip
\subsection*{One-screen summary}

\begin{center}
\renewcommand{\arraystretch}{1.6}
\begin{tabular}{l l l}
\textbf{layer} & \textbf{statement} & \textbf{status} \\
\hline
character & \texttt{traceChi} : concrete $(-1)^{\mathrm{Tr}}$ & \textbf{proved} \\
homogeneity & \texttt{wcorr\_pow\_homog} & \textbf{proved} \\
bridges & \texttt{Delta0\_eq\_Δ}, \texttt{tcount\_eq\_kappa} & \textbf{proved} \\
equivalence & \texttt{correlationBalance\_iff\_count} & \textbf{proved} \\
spectrum & \texttt{kasami\_is\_ab} (3-valued) & \textbf{proved} \\
\rowcolor{leanbg}
\textbf{kernel} & \texttt{kasami\_correlationBalance\_sumzero} & \textbf{open (1 sorry)} \\
count & \texttt{kasami\_triple\_count\_unconditional} & proved \emph{modulo kernel} \\
problem 11 & \texttt{nsucrypto\_problem11} & proved \emph{modulo kernel} \\
\end{tabular}
\end{center}

\bigskip
\noindent
\colorbox{leanbg}{\parbox{\dimexpr\linewidth-2\fboxsep}{\centering
Everything is closed except one number being zero:\\[2pt]
$\displaystyle\sum_{v\neq0}\mathrm{Wcorr}(v v_1)\,\mathrm{Wcorr}(v v_2)\,
\mathrm{Wcorr}\bigl(v(v_1{+}v_2)\bigr)=0.$}}

\bigskip

\section*{Reproducibility}

\noindent
All ``proved'' rows compile in Lean~4 with Mathlib, no \texttt{sorry}, standard
axioms only.  The single open node is
\texttt{kasami\_correlationBalance\_sumzero} (and its general-$m$ companion
\texttt{kasami\_correlationBalance\_sumzero\_m}) in
\texttt{RequestProject/EquivalentContexts/NSUCryptoKernel.lean}.
The assembly and the equivalent-context bridges are in
\texttt{RequestProject/EquivalentContexts/NSUCryptoFoundationalDAG.lean}.

\end{document}
